\subsubsection{Euler's Method}
\label{subsubsec:euler-method}

Euler's method is a straightforward numerical technique for solving differential equations by constructing a continuous set of approximations using tangent lines computed at discrete points.
Given a differential equation with an initial condition:

\[
\frac{dy}{dt} = f(t, y), \quad y(t_0) = y_0,
\]

we start at the point \((t_0, y_0)\) and use the slope to form the tangent line and compute subsequent approximations:

\[
y = y_0 + f(t_0, y_0)(t - t_0)
\quad \rightarrow \quad
y_{n+1} = y_n + f(t_n, y_n) \cdot \Delta t,
\]

where \(\Delta t = t_{n+1} - t_n\) is the fixed step size.

The first equation gives the tangent line approximation from the initial point.
Evaluating this at \(t = t_1\) provides the next point \((t_1, y_1)\).
The iterative process continues using the second equation for all subsequent steps.
As the number of steps increases and \(\Delta t\) decreases, this piecewise linear approximation converges to the true solution \(\phi(t)\).

\subsubsection{Improved Euler's Method}
\label{subsubsec:improved-euler}

The improved Euler's method, also known as Heun's method, is a predictor-corrector approach that achieves better accuracy than the standard Euler method by using an average of slopes.
Using the same initial value problem as in ~\ref{subsubsec:euler-method}, the method proceeds in two sequential steps:

\[
y_{n+1}^* = y_n + f(t_n, y_n) \cdot \Delta t
\quad \rightarrow \quad
y_{n+1} = y_n + \frac{f(t_n, y_n) + f(t_{n+1}, y_{n+1}^*)}{2} \cdot \Delta t.
\]

The first step provides a predictor value using the standard Euler method.
This predicted value is then used to compute a more accurate slope at the next time point.
The second step takes the average of the slopes at the beginning and end of the interval to obtain the corrected value.

By using this averaging technique, the improved Euler method provides a better balance between computational efficiency and accuracy for many ordinary differential equations.

\subsubsection{Runge-Kutta Method}
\label{subsubsec:runge-kutta}

The fourth-order Runge-Kutta method is a widely used numerical technique that provides excellent accuracy while maintaining reasonable computational complexity.
Again, considering the same initial value problem as used in ~\ref{subsubsec:euler-method}, the method requires four slope evaluations per step:

\[
\begin{array}{ll}
k_1 = f(t_n, y_n) &
k_3 = f\left(t_n + \frac{\Delta t}{2}, y_n + \frac{\Delta t}{2}k_2\right) \\[4pt]
k_2 = f\left(t_n + \frac{\Delta t}{2}, y_n + \frac{\Delta t}{2}k_1\right) &
k_4 = f\left(t_n + \Delta t, y_n + \Delta t \cdot k_3\right)
\end{array}
\]

These four slopes are then combined in a weighted average to compute the next value:

\[
y_{n+1} = y_n + \frac{\Delta t}{6}(k_1 + 2k_2 + 2k_3 + k_4).
\]

Despite requiring more function evaluations per step, its higher order of accuracy often allows for larger step sizes while maintaining precision.
The fourth-order Runge-Kutta method strikes an excellent balance between computational cost and accuracy, making it one of the most popular choices for solving ordinary differential equations.